\scanpage{15}
- \mspage{12} -
\[
(6.2)\qquad j=c_4^3/\Delta=1728+c_6^2/\Delta.
\]
\uline{Proposition} 6.3.\quad \uline{Soient} $E$ \uline{et} $E'$ \uline{deux courbes de genre un sur} $S$, \uline{sans fibre additive}.

(I) $\uline{\mathrm{Isom}}_{\inserted{S}}(E,E')$ \uline{est représentable, fini et non ramifié sur} $S$

(II) \uline{si} $j_E=j_F$, \uline{la projection de} $\uline{\mathrm{Isom}}_{\inserted{S}}(E,E')$ \uline{sur} $S$ \uline{est surjective}

(III) \uline{si} $j_E=j_F$ \uline{ne prend pas les valeurs} 0 \uline{et} 1728, $\uline{\mathrm{Isom}}_{\inserted{S}}(E,E')$ \uline{est un revêtement étale de rang} 2 \uline{de} $S$.

Donnons-nous $E$ et $E'$ par des équations (1.3) et (1.3)'. Se donner un isomorphisme entre $E$ et $E'$ revient à se donner $u,r,s$ et $t$, vérifiant les équations (1.6), donc aussi (1.7) et (1.8), $u$ devant de plus être inversible. Par hypothèse, localement sur $S$, $c_4$ ou $\Delta$ est inversible (6.1), et de même $c'_4$ ou $\Delta'$ l'est. Localement sur $S$, les équations (1.8) sont donc des équations de dépendance intégrale pour $u$ et $u^{-1}$, les équations $(1.7)_6$ ou $(1.7)_8$ des équations de dépendance intégrale pour $r$ et enfin $(1.6)_6$ (resp. $(1.6)_2$) est une équation de dépendance intégrale pour $t$ (resp. $s$). Ceci prouve que $\uline{\mathrm{Isom}}(E,E')$ est représentable par un schéma affine fini sur $S$ ; il est même non ramifié car les courbes considérées n'ont pas d'automorphismes infinitésimaux (respectant l'origine).

Lorsque 2 est inversible sur $S$, l'assertion (III) résulte de l'énoncé plus précis.

\uline{Proposition} 6.4.\quad \uline{Soient} $E$ \uline{et} $E'$ \uline{deux courbes de genre un sur} $S$,
\margin{14}
